% Options for packages loaded elsewhere
\PassOptionsToPackage{unicode}{hyperref}
\PassOptionsToPackage{hyphens}{url}
\documentclass[
]{article}
\usepackage{xcolor}
\usepackage{amsmath,amssymb}
\setcounter{secnumdepth}{-\maxdimen} % remove section numbering
\usepackage{iftex}
\ifPDFTeX
  \usepackage[T1]{fontenc}
  \usepackage[utf8]{inputenc}
  \usepackage{textcomp} % provide euro and other symbols
\else % if luatex or xetex
  \usepackage{unicode-math} % this also loads fontspec
  \defaultfontfeatures{Scale=MatchLowercase}
  \defaultfontfeatures[\rmfamily]{Ligatures=TeX,Scale=1}
\fi
\usepackage{lmodern}
\ifPDFTeX\else
  % xetex/luatex font selection
\fi
% Use upquote if available, for straight quotes in verbatim environments
\IfFileExists{upquote.sty}{\usepackage{upquote}}{}
\IfFileExists{microtype.sty}{% use microtype if available
  \usepackage[]{microtype}
  \UseMicrotypeSet[protrusion]{basicmath} % disable protrusion for tt fonts
}{}
\makeatletter
\@ifundefined{KOMAClassName}{% if non-KOMA class
  \IfFileExists{parskip.sty}{%
    \usepackage{parskip}
  }{% else
    \setlength{\parindent}{0pt}
    \setlength{\parskip}{6pt plus 2pt minus 1pt}}
}{% if KOMA class
  \KOMAoptions{parskip=half}}
\makeatother
\setlength{\emergencystretch}{3em} % prevent overfull lines
\providecommand{\tightlist}{%
  \setlength{\itemsep}{0pt}\setlength{\parskip}{0pt}}
\usepackage{bookmark}
\IfFileExists{xurl.sty}{\usepackage{xurl}}{} % add URL line breaks if available
\urlstyle{same}
\hypersetup{
  hidelinks,
  pdfcreator={LaTeX via pandoc}}

\author{}
\date{}

\begin{document}

\section{权利要求书初稿 --- 申请A: 生成式神经表示及训练方法
(G1)}\label{ux6743ux5229ux8981ux6c42ux4e66ux521dux7a3f-ux7533ux8bf7a-ux751fux6210ux5f0fux795eux7ecfux8868ux793aux53caux8badux7ec3ux65b9ux6cd5-g1}

\begin{quote}
日期: 2026-08-11 \textbar{} 状态: 初稿 (交代理人改写为申请格式前) 依据:
专利交底书\_G1\_生成式泛化.md + 权利要求设计报告\_G1.md §3/§4/§5 说明:
文中 \texttt{{[}原则N{]}} 为发明人理论批注, 交代理人前应删除;
各权项技术特征已按六原则组织, 语言全部落在数据结构与流程,
无''可解释/解耦/泛化''抽象词。 关联: M1 宏编译为独立申请B,
见本文档第二部分。
\end{quote}

\begin{center}\rule{0.5\linewidth}{0.5pt}\end{center}

\subsection{第一部分 方法权利要求
(申请A)}\label{ux7b2cux4e00ux90e8ux5206-ux65b9ux6cd5ux6743ux5229ux8981ux6c42-ux7533ux8bf7a}

\textbf{1. 一种生成式神经表示的训练方法, 其特征在于, 包括:}

S1 获取多个概念记录, 每个概念记录包括唯一标识、概念名称以及结构化定义;
所述结构化定义包括定义形式、定义体以及引用列表; 其中,
所述引用列表记录所述定义体引用的全部被引用概念,
所述引用列表中每个被引用概念均已在所述多个概念记录中注册,
从而定义链可枚举、可校验;

S2 根据所述概念记录的排列方法, 将概念递归组合为合法对象;
所述合法对象为可被所述概念体系表达与验证的组合结构;

S3 所述概念记录的定义确定所述概念的求值真值,
即所述概念的运算/判定结果由所述概念记录的定义体决定,
而非由外部规则表决定;

S4 基于所述概念记录生成训练样本集;
所述训练样本集包含所述定义体中出现的概念序列,
使定义链上的概念在训练中得到充分呈现; 所述训练样本集还包括判定样本,
所述判定样本根据所述概念记录的定义确定其真值,
用于使所述神经网络学习概念的判定行为;
属于同一上位概念的多个平行概念并行生成训练样本;

S5 以所述训练样本集训练神经网络, 使所述神经网络同时执行多种概念类别;
其中, 所述概念由相互独立的自由度组合而成,
且所述训练样本集不包含全部自由度组合;

S6 以训练后的神经网络接收包含未在所述训练样本集中出现的自由度组合的输入,
执行得到结果序列, 并验证执行正确性。

\textbf{2. 根据权利要求1所述的方法, 其特征在于:}
所述结构化定义的各字段以数据字段形式承载于所述概念记录;
所述字段的组织存储方式不限定, 包括: 聚合存储于同一条概念记录,
或按字段类别分离存储; 分离存储不改变所述概念记录的表达能力。

\textbf{3. 根据权利要求2所述的方法, 其特征在于:}
所述概念记录按多层投影存储,
所述多层投影包括公设基层、概念层、符号层、排列方法层、呈现层与箭头层;
各层数据经单一注册入口写入,
且同一概念在所述多层投影中各侧记录指向同一唯一标识。

\textbf{4. 根据权利要求1所述的方法, 其特征在于:} 所述定义形式为五类之一,
包括显式定义、归纳定义、递归定义、公理式定义与隐式定义;
其中公理式定义为定义链展开的溯源终点,
定义链展开沿引用列表递归进行并在遇公理式定义时终止。

\textbf{5. 根据权利要求1所述的方法, 其特征在于:} 所述S2中,
所述排列方法为排列方法指示字段; 根据所述排列方法指示,
从语法排列规则集合中读取对应的槽位布局,
所述语法排列规则集合中每条排列规则包括槽位角色词序列;
按所述槽位布局递归填充子表达式, 生成嵌套表达式结构作为所述合法对象;
并经呈现定义对所述嵌套表达式结构执行双向转换验证,
所述呈现定义包括具体记法序列、优先级与结合性,
打印操作与解析操作均由所述呈现定义驱动。

\textbf{6. 根据权利要求1所述的方法, 其特征在于:} 所述S4中,
当任一概念与其相邻概念的搭配对的出现次数达到预设阈值时,
停止为该概念生成更多样本; 所述预设阈值的取值范围为3至10。

\textbf{7. 根据权利要求1所述的方法, 其特征在于:}
将概念之间的有向关系登记为箭头记录,
每条箭头记录包括源概念、目标概念与等价概念; 其中,
所述等价概念为该箭头关系的概念化,
且所述等价概念能够参与更高层级的箭头记录。

\textbf{8. 根据权利要求1所述的方法, 其特征在于:} 所述S3还包括:
根据所述结构化定义自动推导各概念的结构属性,
所述结构属性包括迭代层级、迭代基础、单位元、递归塔规则中的至少一者;
构建统一求值接口,
所述统一求值接口以概念标识与操作数序列为输入、以结果序列为输出,
求值真值由所述概念记录的定义提供;
所述S4基于所述统一求值接口生成训练样本集。

\textbf{9. 根据权利要求1所述的方法, 其特征在于,
所述训练样本集还包括正例与错题负例:}
所述正例为根据所述概念记录的定义确定的正确结果样本;
所述错题负例包括以下至少一者: 与正例同构但结果不同的负例;
针对判定型算符白名单生成的负例; 针对超运算以限定数值域生成的负例。

\textbf{10. 根据权利要求1所述的方法, 其特征在于:} 所述S4中,
属于同一上位概念的多个平行概念并行生成训练样本,
并在训练样本中显式呈现所述平行概念之间的关系。

\textbf{11. 根据权利要求1所述的方法, 其特征在于:}
所述神经网络为单一参数集合的神经网络,
以所述训练样本集对所述单一参数集合执行训练。

\textbf{12. 根据权利要求1所述的方法, 其特征在于:}
所述S5中的训练目标包括位置交叉熵与合法性交叉熵;
所述合法性交叉熵用于区分正例与负例。

\textbf{13. 根据权利要求1所述的方法, 其特征在于:}
所述神经网络的层数小于等于4; 所述训练样本集包含的样本数小于等于10000;
所述神经网络的权重参数总量小于等于4MB。

\textbf{14. 根据权利要求1所述的方法, 其特征在于:}
所述神经网络的层数小于等于2; 所述训练样本集包含的样本数小于等于2000;
所述神经网络的权重参数总量小于等于1MB; 所述S6中,
位置级跨分布准确率大于等于0.99, 且判定口径准确率大于等于0.9。

\textbf{15. 根据权利要求1所述的方法, 其特征在于,
所述S6中的验证进一步包括:} 逐概念泛化率统计,
所述逐概念泛化率为未在训练样本集中出现过的概念自身的预测正确率;
以及带崩概念统计,
用于定位导致判定失败的样本中错误位置真值出现次数最多的概念。

\textbf{16. 根据权利要求1所述的方法, 其特征在于:} 所述S2中,
所述逻辑层、输入层与表达层相互分离;
所述表达层包括优先级、结合性与符号字面量;
所述打印操作将嵌套表达式结构转换为文本,
所述解析操作将文本转换为嵌套表达式结构,
二者均由呈现定义驱动且互为往返验证。

\textbf{17. 根据权利要求8所述的方法, 其特征在于:}
所述S3在推导结构属性时, 同时构建展开执行路径与直接执行路径;
所述展开执行路径沿定义链逐步执行, 所述直接执行路径直接给出执行结果;
所述展开执行路径与所述直接执行路径语义等价,
且所述直接执行路径的计算成本低于所述展开执行路径。

\textbf{18. 根据权利要求1所述的方法, 其特征在于:} 训练与推理过程中,
概念记录的符号标识与概念标识相互分离; 在训练样本集或验证输入中,
至少部分符号标识被置换为其他符号标识时,
所述神经网络对相同概念序列的执行结果保持不变; 其中,
所述符号标识置换不改变所述概念记录的结构化定义。

\textbf{19. 根据权利要求1所述的方法, 其特征在于, 所述训练样本集还包括:}
逻辑命题判定样本,
所述逻辑命题涉及与、或、非、蕴含、当且仅当、异或、与非、或非、同或中的至少一者;
所述逻辑命题判定样本由所述逻辑命题的真值表定义驱动生成。

\textbf{20. 根据权利要求1所述的方法, 其特征在于:} 数值概念采用分层表示,
所述分层包括数值概念层与数位符号层; 其中,
数位符号序列中的每一位按其在序列中的位置确定位权,
所述位权由所述概念记录的定义确定, 并支持不同进制。

\textbf{21. 根据权利要求1所述的方法, 其特征在于,
所述S5中所述训练还包括:} 语法展开, 所述语法展开沿所述概念记录的定义引用,
将每个概念展开为预定义深度的一层或多层被引用概念;
所述展开在训练样本进入神经网络之前执行。

\textbf{22. 根据权利要求1所述的方法, 其特征在于:}
所述神经网络的执行采用自回归方式, 逐位生成结果序列中的元素; 其中,
每次生成元素时, 所述神经网络的权重执行固定次数的前向计算;
所述权重定义状态转移规则, 所述结果序列提供时间维度。

\textbf{23. 根据权利要求1所述的方法, 其特征在于:}
所述概念记录还包括模态标识层;
所述模态标识层将至少一个输入模态标识或输出模态标识映射到所述概念记录;
所述输入模态标识与输出模态标识与所述概念记录的结构化定义解耦;
更换所述输入模态标识, 使所述神经网络接入不同模态的输入;
更换所述输出模态标识, 使所述神经网络产生不同模态的输出或执行动作。

\textbf{24. 根据权利要求1所述的方法, 其特征在于,
所述S6之后还包括反馈步骤:}
将所述神经网络执行得到的结果序列与验证结果转换为概念序列,
作为反馈输入再次送入所述神经网络;
所述神经网络基于所述反馈输入继续执行或修正执行, 直至满足收敛条件; 其中,
反馈环中流通的是可解释的概念序列, 反馈的每个环节可审计、可验证、可解释。

\textbf{25. 根据权利要求1所述的方法, 其特征在于,
所述S6之后还包括多模型协同步骤:} 第一神经网络接收任意模态的输入,
输出第一概念序列; 第二神经网络接收所述第一概念序列作为输入,
以自循环方式迭代执行任务拆解与逐步分析,
将中间分析结果以概念序列形式重新送入自身输入,
迭代收敛后输出结论概念序列; 第三神经网络接收所述结论概念序列,
输出任意模态的内容或执行动作; 其中,
所述第一、第二、第三神经网络基于同一概念记录体系, 网络间以概念序列流转,
无需模态转换层。

\begin{center}\rule{0.5\linewidth}{0.5pt}\end{center}

\subsection{第二部分 装置权利要求
(申请A)}\label{ux7b2cux4e8cux90e8ux5206-ux88c5ux7f6eux6743ux5229ux8981ux6c42-ux7533ux8bf7a}

\textbf{26. 一种生成式神经表示的训练装置, 其特征在于,
包括处理器与存储器;}
所述存储器存储有多个概念记录、语法排列规则集合、呈现定义以及神经网络参数;
所述处理器执行如权利要求1至25中任一项所述的方法。

\textbf{27. 一种计算机可读存储介质, 其上存储有程序, 其特征在于:}
所述程序被处理器执行时实现如权利要求1至25中任一项所述的方法。

\begin{center}\rule{0.5\linewidth}{0.5pt}\end{center}

\subsection{第三部分 M1 宏编译 --- 独立申请B
(分案/分别首申)}\label{ux7b2cux4e09ux90e8ux5206-m1-ux5b8fux7f16ux8bd1-ux72ecux7acbux7533ux8bf7b-ux5206ux6848ux5206ux522bux9996ux7533}

\textbf{1. 一种神经网络计算加速方法, 其特征在于, 包括:}

获取待加速的第一计算操作;

确定所述第一计算操作能够由神经网络已支持的一个或多个基础计算操作组合形成的第一执行路径;

利用所述第一执行路径生成所述第一计算操作对应的多个训练样本;

基于所述多个训练样本更新所述神经网络的至少部分参数,
以建立与所述第一执行路径语义等价的第二执行路径;

对所述第二执行路径与所述第一执行路径的输出结果执行一致性验证;

在验证通过后, 响应于接收到所述第一计算操作, 调用所述第二执行路径执行;
其中,
所述第二执行路径所需的计算步骤数量、序列长度、模型调用次数、存储访问量和/或计算资源消耗低于所述第一执行路径。

\textbf{2. 根据权利要求1所述的方法, 其特征在于:}
所述多个训练样本由所述第一执行路径自动生成, 不依赖外部标注。

\textbf{3. 根据权利要求1所述的方法, 其特征在于:}
所述更新包括以下至少一者:
仅更新嵌入层参数、仅更新适配层参数、仅更新部分网络参数、更新全部网络参数;
其中, 更新过程中保持所述神经网络原有能力的执行正确性。

\textbf{4. 根据权利要求1所述的方法, 其特征在于:}
所述第二执行路径与所述第一执行路径为同一神经网络的两种执行方式,
所述神经网络即作为教师路径又作为学生路径。

\textbf{5. 根据权利要求1所述的方法, 其特征在于, 还包括:}
当所述第二执行路径的执行成本超过预设成本阈值, 或所述一致性验证未通过时,
回退调用所述第一执行路径执行。

\textbf{6. 根据权利要求1所述的方法, 其特征在于:}
所述多个训练样本的数量不超过100个。

\textbf{7. 根据权利要求6所述的方法, 其特征在于:}
所述多个训练样本的数量不超过50个。

\textbf{8. 根据权利要求7所述的方法, 其特征在于:}
所述多个训练样本的数量不超过20个。

\textbf{9. 根据权利要求8所述的方法, 其特征在于:}
所述多个训练样本的数量不超过10个。

\textbf{10. 根据权利要求1所述的方法, 其特征在于:}
所述第一计算操作在更新之前已能够由所述神经网络沿所述第一执行路径正确执行,
且更新过程不扩大所述神经网络的语义闭包。

\textbf{11. 根据权利要求1所述的方法, 其特征在于:}
所述一致性验证采用完整输出序列重建作为判定口径,
所述第二执行路径输出的完整序列与所述第一执行路径输出的完整序列逐位置一致才判定验证通过。

\textbf{12. 根据权利要求1所述的方法, 其特征在于:}
所述第一计算操作包括经由迭代链组合形成的操作,
所述迭代链包括后继、加法、乘法、幂与超运算中的至少两者。

\textbf{13. 根据权利要求1所述的方法, 其特征在于:} 所述更新还包括:
为所述第一计算操作登记新的概念记录与符号标识; 或,
复用所述神经网络已登记的概念记录与符号标识,
不为所述第一计算操作新增登记。

\textbf{14. 根据权利要求1所述的方法, 其特征在于:}
所述多个训练样本由所述第一执行路径自动生成, 或由外部参照求值器生成;
其中, 所述外部参照求值器为与所述神经网络分离的独立求值模块。

\textbf{15. 一种神经网络计算加速装置, 其特征在于, 包括处理器与存储器;}
所述处理器执行如本部分权利要求1至14中任一项所述的方法。

\textbf{16. 一种计算机可读存储介质, 其上存储有程序, 其特征在于:}
所述程序被处理器执行时实现如本部分权利要求1至14中任一项所述的方法。

\begin{center}\rule{0.5\linewidth}{0.5pt}\end{center}

\subsection{附:
交代理人前需确认事项}\label{ux9644-ux4ea4ux4ee3ux7406ux4ebaux524dux9700ux786eux8ba4ux4e8bux9879}

\begin{enumerate}
\def\labelenumi{\arabic{enumi}.}
\tightlist
\item
  删除文中所有 \texttt{{[}原则N{]}} 批注 (本初稿已不含,
  仅在设计报告保留)
\item
  确认 G1 独权 S1-S6 是否需要并入''箭头记录登记''作为前提步骤 (权利要求4
  目前为从属, 可讨论是否并入独权增强防御)
\item
  确认 M1
  独权''利用第一执行路径生成训练样本''是否需补''所述多个训练样本与所述第一执行路径的输出一致''限定
\item
  装置权项中''神经网络参数''与''概念记录''的存储介质分离是否需要
\item
  数据载体权项 (概念投影数据文件结构) 是否单独申请
  (与《数据知识产权》试点重叠, 待代理师判断) --- 见专利申请布局策略.md
  §1 申请C
\item
  效果数据来源: G1 判定口径 0.988 (n=1568, 8/11 05:42 gates6 log), M1 待
  E2 快慢一致性实证
\item
  覆盖缺口补强: A 从属 14 (符号置换不变性)/15 (逻辑命题)/16
  (数值分层)/17 (语法展开)/18 (自回归执行); B 从属 13 (新 token/复用)/14
  (外部 evaluator) --- 见专利申请布局策略.md §4
\item
  \textbf{模态无关扩展 (2026-08-11 发明人确立)}: A 新增从属 19
  (模态标识层, 任意输入/输出模态)/20 (可解释反馈环)/21 (多模型协同);
  独权 S1 已补模态标识层特征; 说明书新增实施例七/八/九 ---
  发明不限于精确逻辑运算系统, 是模态无关的概念执行体系
\item
  \textbf{无代理人自助提交}: 布局/费用/流程/时间线见专利申请布局策略.md
  §5/§6; 占日优先, 授权前资金到位则请代理师改写权利要求
\end{enumerate}

\end{document}
